\documentclass[11pt]{article}
\usepackage[utf8]{inputenc}
\usepackage[T1]{fontenc}
\usepackage{geometry}
\geometry{a4paper, margin=1in}
\usepackage{titlesec}
\usepackage{enumitem}
\usepackage{hyperref}
\usepackage{booktabs}
\usepackage{listings}
\usepackage{xcolor}

\title{\textbf{CredenceAI: Progress Report}}
\author{Anurag Shetye}
\date{\today}

\begin{document}

\maketitle

\section{Project Overview}
CredenceAI is an AI-powered Credit Appraisal Engine specifically designed for the Indian corporate lending landscape. It addresses the ``Data Paradox'' by synthesizing structured and unstructured data through automated ingestion, autonomous research agents, and explainable recommendation engines.

\section{Current Progress}
The system has achieved a functional end-to-end pipeline, from entity initialization to the generation of a comprehensive Credit Appraisal Memo (CAM).

\subsection{Implemented Modules}
\begin{itemize}[noitemsep]
    \item \textbf{Data Ingestion Engine:} Automated parsing of PDF/CSV financial filings (GST returns, ITRs, Bank Statements) using PyMuPDF and GPT-3.5.
    \item \textbf{Autonomous Research Agent:} Real-time web intelligence gathering via Tavily API, focusing on litigation risk, sector headwinds, and promoter track records.
    \item \textbf{Five Cs Scoring Engine:} A hybrid AI-deterministic framework that evaluates Character, Capacity, Capital, Collateral, and Conditions. This module integrates CIBIL-scale (300--900) score conversion via logit transform.
    \item \textbf{CAM Generator:} Synthesis of all data points using GPT-4 into a standard banking-grade Credit Appraisal Memo.
    \item \textbf{GST Reconciliation Hub:} A dashboard for variance analysis between GSTR-1/3B sales and actual bank credits, including circular trading detection.
\end{itemize}

\subsection{Frontend Development}
The UI has been completely overhauled with a professional ``Amber Gold'' institutional banking theme. Key views include:
\begin{itemize}[noitemsep]
    \item Entity Ingestion Dashboard.
    \item Risk Intelligence Visualization (Autonomous Agent).
    \item Five Cs Analysis Radar Matrix.
    \item Interactive Credit Appraisal Memo View.
    \item GST vs. Bank Statement Reconciliation Hub.
\end{itemize}

\section{Accepted Inputs}
The following data sources are currently supported by the ingestion pipeline:
\begin{itemize}[noitemsep]
    \item \textbf{Structured Data:} GST Filings, ITRs, Bank Statements (PDF/CSV).
    \item \textbf{Unstructured Data:} Annual Reports, Financial Statements, Rating Agency Reports.
    \item \textbf{External Data:} News reports, MCA filings, Legal dispute metadata.
    \item \textbf{Primary Data:} Field notes and management interview observations.
\end{itemize}

\section{Remaining and Future Work}
To bridge the intelligence gap further, the following features are slated for development:

\subsection{Immediate Next Steps}
\begin{itemize}[noitemsep]
    \item \textbf{Advanced Document AI:} Integration of OCR + LayoutLM/Donut models for high-fidelity extraction from messy or scanned financial documents.
    \item \textbf{Graph-Based Analysis:} Implementation of Graph Neural Networks (GraphSAGE/GAT) to detect suspicious money loops and circular trading risks.
\end{itemize}

\subsection{Long-Term Roadmap}
\begin{itemize}[noitemsep]
    \item \textbf{News Risk Intelligence Engine:} Deep scraping of regional and sector-specific news with NLP sentiment scoring.
    \item \textbf{Anomaly Detection:} Using Isolation Forests and Autoencoders to flag abnormal revenue patterns in GST and bank flows.
    \item \textbf{Psychometric Analysis:} AI assessment of promoter behavior and management integrity during interviews.
    \item \textbf{Production Deployment:} Cloud infrastructure setup (Docker/AWS) and production-grade database integration.
\end{itemize}

\section{Summary}
CredenceAI has transitioned from a conceptual framework to a working prototype. The core intelligence layers (Ingest, Research, Score) are active and successfully generate actionable banking documents. The upcoming phases will focus on deepening the fraud detection and document extraction capabilities.

\end{document}
